下面给你把文中 **(18)** 和 **(19)** 从 SUPG 弱式一步步推到“单元层面的两行方程”，全程用 LaTeX 表达（符号尽量与论文一致）。

---

## 1. 从强式到 SUPG 试函数

一维平衡方程（论文 Eq. (1)）：
\[
\sigma'(x)+f(x)=0,\qquad \sigma(x)=E(x)\,\varepsilon(x).
\]
写成关于 \(E\) 的形式：
\[
\frac{d}{dx}\big(E(x)\varepsilon(x)\big)=-f(x).
\]

SUPG 的单元内测试函数（论文 Eq. (16)，在单元 \(\Omega^e=[x_i^e,x_{i+1}^e]\)）：
\[
\psi^{e}_a(x)=\phi^{e}_a(x)-\varepsilon(x)\,\tau_e\,\big(\phi^{e}_a\big)'(x),
\qquad a=1,2.
\]
其中 \(\phi_a^e\) 是线性形函数，\(\tau_e\) 为稳定化参数。

---

## 2. SUPG 弱式（单元内积分形式）

由于 \(\psi_a^e\) 含导数项且在单元边界不连续，不能对含 \(\tau\) 的项做全域分部积分；做法是对每个单元分别写加权残值为零：

令残值
\[
R(x):=\sigma'(x)+f(x)=\big(E\varepsilon\big)' + f .
\]
则对每个单元、每个局部测试函数 \(\psi_a^e\)：
\[
\int_{\Omega^e}\psi_a^e \,R \,a(x)\,dx=0.
\]
代入 \(\psi_a^e=\phi_a^e-\varepsilon\tau_e(\phi_a^e)'\)：
\[
\int_{\Omega^e}\phi_a^e\big(\sigma'+f\big)a\,dx
-\int_{\Omega^e}\varepsilon\tau_e(\phi_a^e)'\big(\sigma'+f\big)a\,dx=0.
\]
把它展开成“Galerkin + SUPG”两部分：
\[
\int_{\Omega^e}\phi_a^e \sigma' a\,dx
+\int_{\Omega^e}\phi_a^e f a\,dx
-\int_{\Omega^e}\varepsilon\tau_e(\phi_a^e)'\sigma' a\,dx
-\int_{\Omega^e}\varepsilon\tau_e(\phi_a^e)' f a\,dx=0.
\]

---

## 3. 处理 \(\int \phi_a^e \sigma' a\,dx\)（产生左右端的边界项）

对第一项做分部积分（注意这里是 **\(\phi_a^e\)**，它在单元内连续可分部积分）：
\[
\int_{\Omega^e}\phi_a^e \sigma' a\,dx
=
\Big[\sigma\,a\,\phi_a^e\Big]_{x_i^e}^{x_{i+1}^e}
-\int_{\Omega^e}\sigma\,(a\phi_a^e)'\,dx.
\]
代入 \(\sigma=E\varepsilon\)：
\[
\int_{\Omega^e}\phi_a^e \sigma' a\,dx
=
\Big[\sigma\,a\,\phi_a^e\Big]_{x_i^e}^{x_{i+1}^e}
-\int_{\Omega^e}E\varepsilon\,(a\phi_a^e)'\,dx.
\]

于是单元弱式变为：
\[
-\int_{\Omega^e}E\varepsilon\,(a\phi_a^e)'\,dx
-\int_{\Omega^e}\varepsilon\tau_e(\phi_a^e)'\sigma' a\,dx
+\Big[\sigma a\phi_a^e\Big]_{x_i^e}^{x_{i+1}^e}
+\int_{\Omega^e}f\Big(\phi_a^e-\varepsilon\tau_e(\phi_a^e)'\Big)a\,dx
=0.
\]

把边界项移到右边（更像论文写法）：
\[
\int_{\Omega^e}E\varepsilon\,(a\phi_a^e)'\,dx
+\int_{\Omega^e}\varepsilon\tau_e(\phi_a^e)'\sigma' a\,dx
=
\Big[\sigma a\phi_a^e\Big]_{x_i^e}^{x_{i+1}^e}
+\int_{\Omega^e}f\Big(\phi_a^e-\varepsilon\tau_e(\phi_a^e)'\Big)a\,dx.
\]

---

## 4. 展开 \(\sigma'=(E\varepsilon)'\) 并离散 \(E\)

关键在第二个 SUPG 项：
\[
\int_{\Omega^e}\varepsilon\tau_e(\phi_a^e)'\sigma' a\,dx
=
\int_{\Omega^e}\varepsilon\tau_e(\phi_a^e)'\,(E\varepsilon)' a\,dx.
\]
而
\[
(E\varepsilon)'=E'\varepsilon+E\varepsilon'.
\]

采用与论文一致的线性插值（论文 Eq. (7)）：
\[
E(x)=E_i\,\phi_1^e(x)+E_{i+1}\,\phi_2^e(x).
\]
因此
\[
E'(x)=E_i(\phi_1^e)'(x)+E_{i+1}(\phi_2^e)'(x).
\]

代回 SUPG 项并把未知量 \(E_i,E_{i+1}\) 提出来：

\[
\int_{\Omega^e}\varepsilon\tau_e(\phi_a^e)'\,(E\varepsilon)' a\,dx
=
\int_{\Omega^e}\varepsilon\tau_e(\phi_a^e)'\,
\Big(E'\varepsilon+E\varepsilon'\Big)a\,dx
\]
\[
=
\int_{\Omega^e}\tau_e\,\varepsilon^2(\phi_a^e)'\,E' a\,dx
+\int_{\Omega^e}\tau_e\,\varepsilon\varepsilon'(\phi_a^e)'\,E a\,dx
\]
\[
=
E_i\int_{\Omega^e}\tau_e\varepsilon^2(\phi_a^e)'(\phi_1^e)'a\,dx
+
E_{i+1}\int_{\Omega^e}\tau_e\varepsilon^2(\phi_a^e)'(\phi_2^e)'a\,dx
\]
\[
+
E_i\int_{\Omega^e}\tau_e\varepsilon\varepsilon'(\phi_a^e)'\phi_1^e\,a\,dx
+
E_{i+1}\int_{\Omega^e}\tau_e\varepsilon\varepsilon'(\phi_a^e)'\phi_2^e\,a\,dx.
\]

同样第一项（Galerkin 主项）也展开：
\[
\int_{\Omega^e}E\varepsilon(a\phi_a^e)'\,dx
=
E_i\int_{\Omega^e}\varepsilon(a\phi_a^e)'\phi_1^e\,dx
+
E_{i+1}\int_{\Omega^e}\varepsilon(a\phi_a^e)'\phi_2^e\,dx.
\]

把两部分加起来，得到对每个 \(a=1,2\) 的单元方程：

\[
E_i\int_{\Omega^e}\Big[
\varepsilon(a\phi_a^e)'\phi_1^e
+\tau_e\varepsilon^2(\phi_a^e)'(\phi_1^e)'a
+\tau_e\varepsilon\varepsilon'(\phi_a^e)'\phi_1^e\,a
\Big]\,dx
\]
\[
+
E_{i+1}\int_{\Omega^e}\Big[
\varepsilon(a\phi_a^e)'\phi_2^e
+\tau_e\varepsilon^2(\phi_a^e)'(\phi_2^e)'a
+\tau_e\varepsilon\varepsilon'(\phi_a^e)'\phi_2^e\,a
\Big]\,dx
=
\Big[\sigma a\phi_a^e\Big]_{x_i^e}^{x_{i+1}^e}
+\int_{\Omega^e}f\Big(\phi_a^e-\varepsilon\tau_e(\phi_a^e)'\Big)a\,dx.
\]

---

## 5. 映射到自然坐标，得到论文 Eq. (18) 与 Eq. (19)

自然坐标 \(\xi^e\in[0,1]\)（论文 Eq. (8)）：
\[
x(\xi^e)=x_i^e(1-\xi^e)+x_{i+1}^e\xi^e,\qquad
dx=J^e\,d\xi^e,\quad J^e=\frac{dx}{d\xi^e}.
\]
局部形函数（论文 Eq. (6)）：
\[
\phi_1^e=1-\xi^e,\qquad \phi_2^e=\xi^e,
\qquad (\phi_a^e)'=\frac{1}{J^e}\frac{d\phi_a^e}{d\xi^e}.
\]

把上面的单元方程写成在 \([0,1]\) 的积分：

对 \(a=1\)（即测试函数选 \(\phi_1^e\)），边界项
\[
\Big[\sigma a\phi_1^e\Big]_{0}^{1}
=\sigma(\xi^e=1)a(\xi^e=1)\underbrace{\phi_1^e(1)}_{0}
-\sigma(\xi^e=0)a(\xi^e=0)\underbrace{\phi_1^e(0)}_{1}
=-\sigma(\xi^e=0)a(\xi^e=0).
\]
于是得到（这就是论文 Eq. (18) 的结构）：
\[
E_i\int_0^1
\Big[
\varepsilon (a\phi_1^e)'\phi_1^e
+\tau_e\varepsilon^2(\phi_1^e)'(\phi_1^e)'a
+\tau_e\varepsilon\varepsilon'(\phi_1^e)'\phi_1^e\,a
\Big]J^e\,d\xi^e
\]
\[
+
E_{i+1}\int_0^1
\Big[
\varepsilon (a\phi_1^e)'\phi_2^e
+\tau_e\varepsilon^2(\phi_1^e)'(\phi_2^e)'a
+\tau_e\varepsilon\varepsilon'(\phi_1^e)'\phi_2^e\,a
\Big]J^e\,d\xi^e
\]
\[
=
-\sigma(\xi^e=0)a(\xi^e=0)
+\int_0^1 f\Big(\phi_1^e-\varepsilon\tau_e(\phi_1^e)'\Big)a\,J^e\,d\xi^e.
\tag{18}
\]

对 \(a=2\)，边界项
\[
\Big[\sigma a\phi_2^e\Big]_{0}^{1}
=\sigma(1)a(1)\underbrace{\phi_2^e(1)}_{1}
-\sigma(0)a(0)\underbrace{\phi_2^e(0)}_{0}
=+\sigma(\xi^e=1)a(\xi^e=1).
\]
于是得到（论文 Eq. (19)）：
\[
E_i\int_0^1
\Big[
\varepsilon (a\phi_2^e)'\phi_1^e
+\tau_e\varepsilon^2(\phi_2^e)'(\phi_1^e)'a
+\tau_e\varepsilon\varepsilon'(\phi_2^e)'\phi_1^e\,a
\Big]J^e\,d\xi^e
\]
\[
+
E_{i+1}\int_0^1
\Big[
\varepsilon (a\phi_2^e)'\phi_2^e
+\tau_e\varepsilon^2(\phi_2^e)'(\phi_2^e)'a
+\tau_e\varepsilon\varepsilon'(\phi_2^e)'\phi_2^e\,a
\Big]J^e\,d\xi^e
\]
\[
=
+\sigma(\xi^e=1)a(\xi^e=1)
+\int_0^1 f\Big(\phi_2^e-\varepsilon\tau_e(\phi_2^e)'\Big)a\,J^e\,d\xi^e.
\tag{19}
\]

---

如果你把你“文中 Eq.(17)”那一行原始公式（带求和那行）也贴成清晰的 LaTeX/截图，我可以把 **(17) → (18)(19)** 的中间步骤写得更“逐行对照论文”，包括每一项从哪来的、符号为什么是正/负。